\documentclass[11pt,letterpaper]{article}
\usepackage[T1]{fontenc}
\usepackage{lmodern}
\usepackage[margin=1in,headheight=14pt]{geometry}
\usepackage{amsmath,amssymb,amsthm,mathtools}
\usepackage{microtype,booktabs,longtable,array}
\usepackage[dvipsnames]{xcolor}
\usepackage{hyperref}
\usepackage{fancyhdr}
\hypersetup{colorlinks=true,linkcolor=MidnightBlue,citecolor=MidnightBlue,urlcolor=MidnightBlue,
 pdftitle={Balanced Stieltjes Moments: Bernoulli Closure, Harmonic Filters, and Polylogarithmic Resonances},
 pdfauthor={Research continuation for the ProveIt project}}
\pagestyle{fancy}\fancyhf{}
\fancyhead[L]{\small Balanced Stieltjes moments}
\fancyhead[R]{\small ProveIt research continuation}
\fancyfoot[C]{\thepage}
\setlength{\parskip}{4pt}
\setlength{\parindent}{1.25em}
\setcounter{tocdepth}{2}
\numberwithin{equation}{section}
\newtheorem{theorem}{Theorem}[section]
\newtheorem{lemma}[theorem]{Lemma}
\newtheorem{proposition}[theorem]{Proposition}
\newtheorem{corollary}[theorem]{Corollary}
\theoremstyle{definition}\newtheorem{definition}[theorem]{Definition}
\theoremstyle{remark}\newtheorem{remark}[theorem]{Remark}
\DeclareMathOperator{\Li}{Li}
\DeclareMathOperator{\Ree}{Re}
\DeclareMathOperator{\FP}{FP}
\newcommand{\dd}{\,\mathrm{d}}
\newcommand{\ii}{\mathrm{i}}
\newcommand{\Ccal}{\mathcal C}
\newcommand{\Acal}{\mathcal A}
\newcommand{\Jcal}{\mathcal J}
\newcommand{\Lcal}{\mathcal L}
\newcommand{\Kcal}{\mathcal K}
\newcommand{\Qcal}{Q}
\newcommand{\eps}{\varepsilon}
\newcommand{\Z}{\mathbb Z}
\newcommand{\N}{\mathbb N}
\newcommand{\C}{\mathbb C}
\newcommand{\R}{\mathbb R}
\newcommand{\rf}[2]{(#1)_{#2}}
\allowdisplaybreaks[2]

\title{\textbf{Balanced Stieltjes Moments}\\[5pt]
\large Bernoulli Closure, Harmonic Filters,\\and Polylogarithmic Resonances}
\author{A research continuation prepared for Vladimir Reshetnikov\\
\small The ProveIt project}
\date{11 October 2026}
\begin{document}
\maketitle
\begin{abstract}
We construct a binomially balanced family of generalized Stieltjes functions
whose first several inverse-power tails cancel identically. Its Mellin
transform is a product of a Hurwitz zeta function, the reflection factor
$\pi\csc(\pi u)$, and one Gamma-quotient coefficient. Every convergent
positive-integer moment consequently reduces to a Bernoulli polynomial
times a finite rational combination of ordinary zeta values. At a zero
of the relevant Bernoulli polynomial, one additional moment converges;
its value is instead a Hurwitz-zeta order derivative, or equivalently an
explicit combination of polylogarithm order derivatives.

We prove uniqueness of the normalized homogeneous balancing, a necessary
and sufficient finite-shift cancellation criterion, and a two-variable
formula for all logarithmic moments of filtered Stieltjes constants.
The last formula has coefficients generated by finite harmonic sums and
ordinary zeta values. It supplies an exact distinction between elementary
and complete homogeneous harmonic coordinates. Shift derivatives and
anchored primitives give further convergent identities, including an
all-order centered vanishing family and an odd-zeta evaluation at the
first exceptional centered moment.

The classical beta integral, Euler--Maclaurin formula, and Hurwitz
functional equation are credited inputs. The specific all-order
organization is offered as a proved research contribution for independent
review, not as a certified literature-priority claim. Existing Gaussian
$S_6$ and revised $S_8$ conjectures are not resolved here.
\end{abstract}
\vspace{4pt}
\noindent\textbf{Keywords.} Generalized Stieltjes constants; Mellin transforms;
Bernoulli polynomials; Gamma quotients; harmonic sums; polylogarithm order
derivatives; ordinary convergent integrals.
\clearpage
{\setlength{\parskip}{0pt}\small\tableofcontents}
\clearpage

\section{Scope, motivation, and principal results}\label{bsm:sec:scope}

The ProveIt manuscript develops bridges among polylogarithms, Hurwitz
jets, Stieltjes calculus, and Gamma correlations. The inspected
harmonic-generating notes explicitly distinguish their known constant
terms from genuinely mixed harmonic--Stieltjes coefficients
\cite{RepoHarmonic,RepoQuestions}. The inspected Mellin continuation
integrates polylogarithms in their argument and obtains Hurwitz and
Gamma expressions \cite{RepoMellin}. We study a different integral:
a Mellin transform in the \emph{positive translation variable} $x$ of
Stieltjes functions at $x+a$.

A term such as $\gamma_q(x+a)$ by itself has a growing logarithmic
asymptotic expansion, so it has no positive-integer moment on
$(0,\infty)$. The combination
\begin{equation}\label{bsm:eq:C-intro}
 \Ccal_q(x,a)=\frac{\log^{q+1}x}{q+1}
 +\sum_{j=0}^{q}(-1)^j\binom qj\gamma_j(x+a)\log^{q-j}x
\end{equation}
removes not merely that growth but all inverse powers $x^{-1},\ldots,x^{-q}$.
The resulting ordinary integral satisfies
\begin{equation}\label{bsm:eq:intro-master}
 \int_0^\infty x^{u-1}\Ccal_q(x,a)\dd x
 =\frac{\pi}{\sin\pi u}\,\zeta(1-u,a)\Qcal_q(u),
 \qquad 0<\Ree u<q+1,
\end{equation}
where
\begin{equation}\label{bsm:eq:intro-Q}
 \sum_{q\ge0}\frac{\Qcal_q(u)}{q!}z^q
 =\frac{\Gamma(u+z)}{\Gamma(u)\Gamma(1+z)}.
\end{equation}
Interior integer singularities on the right are removable. Their
removal is an essential part of the theorem, not an informal rule for
multiplying zero by infinity.

For $q\ge m\ge1$, let
\begin{equation}\label{bsm:eq:intro-kappa}
 \kappa_{m,q}=\frac{(-1)^{m+q}q!}{m}
 \sum_{j=0}^{m-1}(-1)^j e_j(1,1/2,\ldots,1/(m-1))\zeta(q+1-j).
\end{equation}
Here $e_0=1$ and an empty positive-degree elementary symmetric sum is
zero. The complete ordinary moment is
\begin{equation}\label{bsm:eq:intro-integer}
 \boxed{\displaystyle
 \int_0^\infty x^{m-1}\Ccal_q(x,a)\dd x
       =\kappa_{m,q}B_m(a).}
\end{equation}
For example,
\begin{align}
 \int_0^\infty\Ccal_1(x,a)\dd x
   &=(a-\tfrac12)\zeta(2),\label{bsm:eq:intro-first}\\
 \int_0^\infty x\Ccal_2(x,a)\dd x
   &=\bigl(\zeta(3)-\zeta(2)\bigr)
                  (a^2-a+\tfrac16).\label{bsm:eq:intro-second}
\end{align}
A distinct cancellation occurs at a Bernoulli root:
\begin{equation}\label{bsm:eq:intro-exception}
 B_m(a)=0\quad\Longrightarrow\quad
 \int_0^\infty x^{m-1}\Ccal_{m-1}(x,a)\dd x
           =(-1)^{m-1}\zeta^{[1]}(1-m,a).
\end{equation}
In particular,
\begin{equation}\label{bsm:eq:intro-odd}
 \boxed{\displaystyle
 \int_0^\infty x^2\Ccal_2(x,\tfrac12)\dd x
             =\frac{3\zeta(3)}{16\pi^2}.}
\end{equation}
Superscript $[r]$ always denotes differentiation in the \emph{order}
of zeta or polylogarithm, whereas parentheses on $\psi^{(r)}$ and
$\gamma_n^{(r)}$ denote differentiation in their argument.

The article then classifies which finite combinations of shifts admit
additional moments and gives a full mixed logarithmic formula. The
classification concerns exact analytic convergence and reduction in
specified coordinates. It is not an arithmetic-independence theorem.

\subsection{What is and is not claimed}

The beta transform underlying \eqref{bsm:eq:intro-master}, the Gamma
expansions in \eqref{bsm:eq:intro-Q}, Bernoulli special values, and the
Hurwitz Fourier formula are classical \cite{DLMFHurwitz,DLMFGamma,Zagier}.
Our contribution is the combined cancellation theorem, its uniqueness
within a precisely stated ansatz, the two different resonance mechanisms,
and the finite-shift and mixed-moment consequences proved below. A
focused source comparison did not establish an earlier identical
all-order formulation. That observation is not an exhaustive priority
search, particularly for the unexpanded incoming archives.

No theorem below evaluates the unresolved coefficients $R_{r,k}(a)$ of
the repository's harmonic Barnes--Hurwitz family. Our $\gamma_n(x+a)$
are ordinary generalized Stieltjes constants, not those higher harmonic
constants. No result is asserted for the Gaussian $S_6$ or revised
$S_8$ period residuals. These distinctions avoid silently changing the
problem whose solution is claimed.

\section{Balanced Stieltjes functions and cancellation}\label{bsm:sec:balanced}

\subsection{Conventions and the holomorphic generator}
Throughout the main theorems, $a>0$, $x>0$, and $\log x$ is real.
The normalization is
\begin{equation}\label{bsm:eq:stieltjes}
 \zeta(1+z,a)=\frac1z+
             \sum_{n=0}^\infty\frac{(-1)^n\gamma_n(a)}{n!}z^n,
 \qquad \gamma_0(a)=-\psi(a).
\end{equation}
Bernoulli polynomials have the standard normalization~\cite{DLMFBernoulli}
and generating function
\begin{equation}\label{bsm:eq:bernoulli-generator}
 \frac{t e^{at}}{e^t-1}=\sum_{k=0}^\infty B_k(a)\frac{t^k}{k!},
 \qquad B_1(a)=a-\frac12.
\end{equation}
Put $P_n(z)=(1+z)_n$, with $P_0=1$. Define
\begin{equation}\label{bsm:eq:F}
 F(z;x,a)=x^z\zeta(1+z,x+a)-\frac1z.
\end{equation}
The singularity at $z=0$ is removable. In fact $F$ is entire in $z$
for each $x,a>0$, since the only order pole of Hurwitz zeta has been
removed. Direct multiplication of \eqref{bsm:eq:stieltjes} gives
\begin{equation}\label{bsm:eq:C-generator}
 \Ccal_q(x,a)=\left.\partial_z^qF(z;x,a)\right|_{z=0},
 \qquad F(z;x,a)=\sum_{q\ge0}\Ccal_q(x,a)\frac{z^q}{q!}.
\end{equation}
The first functions are
\begin{align}
 \Ccal_0(x,a)&=\log x-\psi(x+a),\label{bsm:eq:C0}\\
 \Ccal_1(x,a)&=-\psi(x+a)\log x-\gamma_1(x+a)
                                      +\tfrac12\log^2x,\label{bsm:eq:C1}\\
 \Ccal_2(x,a)&=-\psi(x+a)\log^2x-2\gamma_1(x+a)\log x
                          +\gamma_2(x+a)+\tfrac13\log^3x.\label{bsm:eq:C2}
\end{align}

\begin{lemma}[Normally uniform inverse-power expansion]\label{bsm:lem:tail}
For every integer $K\ge1$, uniformly for $z$ in a sufficiently small
closed disc about zero and for $a$ in a compact subset of $(0,\infty)$,
\begin{equation}\label{bsm:eq:F-tail}
 F(z;x,a)=\sum_{k=1}^{K}
       \frac{(-1)^k B_k(a)}{k!}P_{k-1}(z)x^{-k}
                       +O(x^{-K-1})\qquad(x\longrightarrow\infty).
\end{equation}
Consequently, for fixed $q\ge0$,
\begin{equation}\label{bsm:eq:C-tail}
 \Ccal_q(x,a)=\sum_{k=q+1}^{K}
       \frac{(-1)^k B_k(a)}{k!}P_{k-1}^{(q)}(0)x^{-k}
                       +O(x^{-K-1}),
\end{equation}
where the sum is empty when $K\le q$. In particular,
\begin{equation}\label{bsm:eq:leading-tail}
 \Ccal_q(x,a)=\frac{(-1)^{q+1}B_{q+1}(a)}{q+1}x^{-q-1}
                              +O(x^{-q-2}).
\end{equation}
At zero, $\Ccal_q(x,a)=O(1+|\log x|^{q+1})$, uniformly on the same
compact sets of $a$.
\end{lemma}

\begin{proof}
Euler--Maclaurin applied to the defining Hurwitz series, followed by
expansion in $a/x$, gives the standard shifted asymptotic formula
\[
 \zeta(s,x+a)=\frac{x^{1-s}}{s-1}
 +\sum_{k=1}^{K}\frac{(-1)^k B_k(a)}{k!}(s)_{k-1}x^{1-s-k}
 +O\bigl(|x^{1-s}|x^{-K-1}\bigr).
\]
It is enough to use more Euler--Maclaurin terms than $K$; their bounded
periodic-Bernoulli remainder yields this estimate uniformly for $s$
in a compact set near $1$ after the leading pole is removed. These are
the usual local uniformity properties of the Hurwitz asymptotic
expansion \cite[\S25.11(iii),(xii)]{DLMFHurwitz}.

For completeness, the Bernoulli coefficient can be checked without
assuming the shifted form. The leading integral, half-endpoint term,
and the $r$th even Euler--Maclaurin term contribute at $x^{-k}$,
after multiplication by $x^{s-1}$, a common factor $(s)_{k-1}$ times
\[
 \frac{(-1)^k a^k}{k!}
 +\frac{(-1)^{k-1}a^{k-1}}{2(k-1)!}
 +\sum_{r=1}^{\lfloor k/2\rfloor}
   \frac{B_{2r}(-1)^{k-2r}a^{k-2r}}{(2r)!(k-2r)!}.
\]
By \eqref{bsm:eq:bernoulli-generator}, this is
$(-1)^kB_k(a)/k!$. Set $s=1+z$ to obtain \eqref{bsm:eq:F-tail}.
Uniformity on a slightly larger $z$-disc permits Cauchy differentiation
of its remainder, with the same inverse-power bound and no additional
logarithm of $x$.

Since $P_{k-1}$ has degree $k-1$, its $q$th derivative vanishes for
$k\le q$. Also $P_q^{(q)}(0)=q!$, giving
\eqref{bsm:eq:leading-tail}. Finally, for $x$ near zero all $\gamma_j(x+a)$
are uniformly bounded on compact positive $a$-intervals. Formula
\eqref{bsm:eq:C-intro} proves the last estimate.
\end{proof}

\begin{corollary}[Ordinary convergence strip]\label{bsm:cor:strip}
For each $q\ge0$ the integral
\begin{equation}\label{bsm:eq:J-def}
 \Jcal_q(u,a)=\int_0^\infty x^{u-1}\Ccal_q(x,a)\dd x
\end{equation}
converges absolutely and is holomorphic in $u$ on
$0<\Ree u<q+1$. It is locally uniformly convergent in $a>0$.
If $B_{q+1}(a)=0$, it converges on the larger strip
$0<\Ree u<q+2$ for that fixed $a$.
\end{corollary}
\begin{proof}
Apply Lemma~\ref{bsm:lem:tail} at the two endpoints. On a compact substrip,
additional powers of $\log x$ remain integrable, proving holomorphy.
\end{proof}

\subsection{Why the binomial coefficients are forced}

\begin{theorem}[Uniqueness in the homogeneous logarithmic ansatz]\label{bsm:thm:unique}
Fix $q\ge0$. Consider constants $d_0,\ldots,d_q,b$ with $d_0=1$ and
\begin{equation}\label{bsm:eq:ansatz}
 D_q(x,a)=b\log^{q+1}x+
               \sum_{j=0}^q d_j\gamma_j(x+a)\log^{q-j}x.
\end{equation}
Suppose $D_q(x,a)=O(x^{-q-1})$ as $x\to\infty$ for every $a$ in some
nonempty open positive interval. Then
\begin{equation}\label{bsm:eq:unique-coeff}
 d_j=(-1)^j\binom qj,
 \qquad b=\frac1{q+1},
 \qquad D_q=\Ccal_q.
\end{equation}
\end{theorem}

\begin{proof}
For $q=0$, cancellation of the leading logarithm forces $b=1$.
Let $q\ge1$. Write $L=\log x$ and $d(t)=\sum_{j=0}^q d_jt^j$.
The expansion underlying Lemma~\ref{bsm:lem:tail} gives, at each inverse
power $x^{-k}$, the coefficient in \eqref{bsm:eq:ansatz}
\[
 \frac{(-1)^kB_k(a)}{k!}
 \sum_{r=0}^{\min(q,k-1)}(-1)^rP_{k-1}^{(r)}(0)L^{q-r}
                  \sum_{j=r}^{q}\binom jr d_j.
\]
For $1\le k\le q$ this polynomial in $L$ must vanish. A nonzero
Bernoulli polynomial cannot vanish throughout the prescribed open
$a$-interval. Choosing $k=r+1$ for $r=0,\ldots,q-1$ and using
$P_r^{(r)}(0)=r!$ therefore forces $d^{(r)}(1)=0$.
Thus $d(t)$ is a constant multiple of $(1-t)^q$; $d(0)=1$ fixes the
constant. The remaining leading logarithm vanishes precisely when
\[
 b=\sum_{j=0}^q\frac{(-1)^j\binom qj}{j+1}
   =\int_0^1(1-t)^q\dd t=\frac1{q+1}.
\]
Conversely these coefficients have the required decay by
Lemma~\ref{bsm:lem:tail}.
\end{proof}

This uniqueness statement is restricted to \eqref{bsm:eq:ansatz}. It does
not classify arbitrary combinations with rational $x$-dependent
coefficients or arbitrary differential operators.

\subsection{Exact shift and averaging relations}
The elementary Hurwitz recurrence gives
\begin{equation}\label{bsm:eq:shift}
 \boxed{\displaystyle
 \Ccal_q(x,a+1)-\Ccal_q(x,a)
       =-\frac{\log^q(x/(x+a))}{x+a}.}
\end{equation}
This is an identity of ordinary functions, including $q=0$.
Integrating $\partial_a\zeta(z,x+a)=-z\zeta(1+z,x+a)$ over
$0\le a\le1$ instead gives, for each fixed $x>0$,
\begin{equation}\label{bsm:eq:mean}
 \int_0^1F(z;x,a)\dd a=0,
 \qquad \int_0^1\Ccal_q(x,a)\dd a=0.
\end{equation}
The apparent $z=0$ singularity is removed before differentiation.
Equation~\eqref{bsm:eq:mean} concerns a finite shift interval; it does not
assert an unproved interchange of two unbounded integrals.

\section{The Mellin theorem and its analytic proof}\label{bsm:sec:mellin}

\subsection{The harmonic coefficient}
For $\Ree u>0$, define $\Qcal_q(u)$ by \eqref{bsm:eq:intro-Q} near $z=0$.
Set
\begin{align}
 H_{u-1}^{(1)}&=\psi(u)+\gamma,\notag\\
 H_{u-1}^{(r)}&=\zeta(r)-\zeta(r,u)\qquad(r\ge2).
 \label{bsm:eq:general-harmonic}
\end{align}
These continue the usual generalized harmonic numbers. The Gamma
Taylor formula gives
\begin{equation}\label{bsm:eq:Q-log}
 \frac{\Gamma(u+z)}{\Gamma(u)\Gamma(1+z)}
 =\exp\left(\sum_{r\ge1}\frac{(-1)^{r-1}}{r}
                              H_{u-1}^{(r)}z^r\right).
\end{equation}
Consequently
\begin{align}
 \Qcal_0&=1,&\Qcal_1&=H^{(1)},\notag\\
 \Qcal_2&=(H^{(1)})^2-H^{(2)},&
 \Qcal_3&=(H^{(1)})^3-3H^{(1)}H^{(2)}+2H^{(3)}.
 \label{bsm:eq:Q-low}
\end{align}
At a positive integer $m$,
\begin{equation}\label{bsm:eq:Q-integer}
 \frac{\Gamma(m+z)}{\Gamma(m)\Gamma(1+z)}
     =E_m(z):=\prod_{j=1}^{m-1}\left(1+\frac zj\right),
 \qquad \Qcal_q(m)=q!e_q^{(m)},
\end{equation}
where $e_q^{(m)}=[z^q]E_m(z)$ and $e_q^{(m)}=0$ for $q\ge m$.
This is an \emph{elementary}, not complete homogeneous, harmonic
polynomial. The distinction will reverse in a different direction of
the mixed generating function in Section~\ref{bsm:sec:mixed}.

\begin{theorem}[Balanced Mellin transform]\label{bsm:thm:master}
For $a>0$, $q\ge0$, and $0<\Ree u<q+1$,
\begin{align}\label{bsm:eq:master}
 \Jcal_q(u,a)
  &=\Gamma(1-u)\zeta(1-u,a)
     \left.\partial_z^q\frac{\Gamma(u+z)}{\Gamma(1+z)}\right|_{z=0}\notag\\
  &=\frac{\pi}{\sin\pi u}\,\zeta(1-u,a)\Qcal_q(u).
\end{align}
At $u=1,\ldots,q$, the entire right-hand expression is evaluated by its
removable limit. These are values of absolutely convergent integrals.
\end{theorem}

\begin{proof}
First take $\Ree p>0$ and $\Ree v>\Ree p+1$. Termwise beta integration,
justified by absolute convergence, gives
\begin{equation}\label{bsm:eq:raw-beta}
 \int_0^\infty x^{p-1}\zeta(v,x+a)\dd x
 =\frac{\Gamma(p)\Gamma(v-p)}{\Gamma(v)}\zeta(v-p,a).
\end{equation}
Indeed the integral of $(x+n+a)^{-v}$ is
$(n+a)^{p-v}B(p,v-p)$.

We need the adjacent strip, not an unjustified substitution into this
integral. Split at $x=1$ and subtract only the leading tail there. The
meromorphic continuation of the left side of \eqref{bsm:eq:raw-beta} is
\begin{align}\label{bsm:eq:cutoff-continuation}
 &\int_0^1x^{p-1}\zeta(v,x+a)\dd x\notag\\
 &\quad+\int_1^\infty x^{p-1}
      \left(\zeta(v,x+a)-\frac{x^{1-v}}{v-1}\right)\dd x
       +\frac1{(v-1)(v-p-1)}.
\end{align}
The subtracted tail converges for $\Ree p<\Ree v$. In the overlapping
strip
\[
 \Ree p>0,\qquad \Ree(v-1)<\Ree p<\Ree v,
\]
the last term in \eqref{bsm:eq:cutoff-continuation} is the negative of
$\int_0^1x^{p-v}\dd x/(v-1)$. Thus the continued expression is the
\emph{ordinary} whole-axis difference integral
\begin{equation}\label{bsm:eq:subtracted-beta}
 \int_0^\infty x^{p-1}
      \left(\zeta(v,x+a)-\frac{x^{1-v}}{v-1}\right)\dd x
 =\frac{\Gamma(p)\Gamma(v-p)}{\Gamma(v)}\zeta(v-p,a).
\end{equation}
This cutoff argument is a standard Mellin-continuation principle
\cite{Zagier}; it specifies all terms and does not assign a value to
the separate divergent integral of a monomial.

Now put $v=1+z$ and $p=u+z$. For $0<\Ree u<1$ and $z$ sufficiently
near zero, the left side becomes
\[
 \int_0^\infty x^{u-1}F(z;x,a)\dd x,
\]
and the right side becomes
\begin{equation}\label{bsm:eq:master-generator}
 \frac{\Gamma(u+z)\Gamma(1-u)}{\Gamma(1+z)}\zeta(1-u,a).
\end{equation}
Near infinity, Lemma~\ref{bsm:lem:tail} is uniform in $z$. Near zero,
choose $|\Ree z|<\Ree u/2$; the regularized leading term
$(x^z-1)/z$ is bounded by a constant times
$(1+|\log x|)x^{-|\Ree z|}$, and the remaining local Hurwitz family
has the same integrable type. Cauchy differentiation or dominated
parameter differentiation now gives the first line of
\eqref{bsm:eq:master} on $0<\Ree u<1$.

Corollary~\ref{bsm:cor:strip} extends the left side holomorphically to
$0<\Ree u<q+1$. The right side is meromorphic there. Analytic
continuation makes the two equal and forces every pole in this strip
to be removable. Finally, Euler's reflection formula
$\Gamma(u)\Gamma(1-u)=\pi/\sin\pi u$ gives the second line.
\end{proof}

\begin{remark}[Why the differentiated strip is larger]
The undifferentiated integral of $F$ generally has only the strip
$0<\Ree u<1$. Differentiating it there is legitimate. The \emph{resulting}
function $\Ccal_q$ has a larger convergence strip because derivatives
annihilate finite-degree Pochhammer polynomials. One must not differentiate
a divergent integral at an arbitrary integer $u$ and call that the proof.
\end{remark}

\section{All positive-integer moments}\label{bsm:sec:integer}

\begin{theorem}[Bernoulli closure]\label{bsm:thm:integer}
For $a>0$ and integers $q\ge m\ge1$, the coefficient
$\kappa_{m,q}$ in \eqref{bsm:eq:intro-kappa} is finite and nonzero, and
\begin{equation}\label{bsm:eq:integer-main}
 \Jcal_q(m,a)=\kappa_{m,q}B_m(a).
\end{equation}
Equivalently,
\begin{equation}\label{bsm:eq:kappa-positive}
 \kappa_{m,q}=
 \frac{(-1)^{q+1}q!}{m}
 \sum_{n=m}^\infty\frac{\binom{n-1}{m-1}}{n^{q+1}}.
\end{equation}
The infinite series here is a second representation of a finite zeta
combination, not an additional unresolved constant.
\end{theorem}

\begin{proof}
At $u=m$, \eqref{bsm:eq:Q-integer} gives $\Qcal_q(m)=0$. Since
$\pi\csc\pi u=(-1)^m/(u-m)+O(u-m)$, Theorem~\ref{bsm:thm:master} gives
\begin{equation}\label{bsm:eq:integer-limit}
 \Jcal_q(m,a)=\frac{(-1)^{m+1}}m B_m(a)\Qcal_q'(m).
\end{equation}
We used the classical value $\zeta(1-m,a)=-B_m(a)/m$.
Differentiating the Gamma quotient in $u$ gives
\begin{equation}\label{bsm:eq:Q-derivative}
 \Qcal_q'(m)=q![z^q]E_m(z)\bigl(\psi(m+z)-\psi(m)\bigr).
\end{equation}
The digamma recurrence writes the last factor as
\[
 \psi(1+z)-\psi(1)
       +\sum_{j=1}^{m-1}\left(\frac1{j+z}-\frac1j\right).
\]
After multiplication by $E_m$, the finite rational sum is a polynomial
of degree at most $m-1$. It cannot contribute at degree $q\ge m$.
The Taylor series
\[
 \psi(1+z)-\psi(1)=\sum_{r\ge1}(-1)^{r+1}\zeta(r+1)z^r
\]
therefore gives
\begin{equation}\label{bsm:eq:Qprime}
 \Qcal_q'(m)=q!\sum_{j=0}^{m-1}
                    e_j^{(m)}(-1)^{q-j+1}\zeta(q-j+1).
\end{equation}
Substitution into \eqref{bsm:eq:integer-limit} proves
\eqref{bsm:eq:integer-main} with exactly the sign in
\eqref{bsm:eq:intro-kappa}.

To obtain \eqref{bsm:eq:kappa-positive}, expand every zeta value into its
absolutely convergent series. The coefficient at $n$ is
\[
 \sum_{j=0}^{m-1}(-1)^j e_j^{(m)}n^j
 =\prod_{j=1}^{m-1}(1-n/j)
 =(-1)^{m-1}\binom{n-1}{m-1}.
\]
It vanishes at $n=1,\ldots,m-1$. The remaining series converges since
$q\ge m$, and all its terms are positive. Thus $\kappa_{m,q}\ne0$.
\end{proof}

\subsection{Explicit evaluations and a centered vanishing family}
For every $q\ge1$,
\begin{equation}\label{bsm:eq:unweighted}
 \boxed{\displaystyle
 \int_0^\infty\Ccal_q(x,a)\dd x
       =(-1)^{q+1}q!(a-\tfrac12)\zeta(q+1).}
\end{equation}
In particular,
\begin{equation}\label{bsm:eq:centered-zero}
 \int_0^\infty\left[
 \frac{\log^{q+1}x}{q+1}
 +\sum_{j=0}^q(-1)^j\binom qj
                \gamma_j(x+\tfrac12)\log^{q-j}x\right]\dd x=0
 \quad(q\ge1).
\end{equation}
Each complete bracket is integrated as one function. Its summands
need not have separate convergent integrals.

At $q=1$, $a=1$, this gives
\begin{equation}\label{bsm:eq:pi12}
 \int_0^\infty\left[-\psi(x+1)\log x-\gamma_1(x+1)
                         +\frac12\log^2x\right]\dd x=\frac{\pi^2}{12}.
\end{equation}
At $q=2$ the first weighted evaluation is
\begin{align}\label{bsm:eq:weighted-C2}
 &\int_0^\infty x\left[-\psi(x+a)\log^2x
       -2\gamma_1(x+a)\log x+\gamma_2(x+a)
       +\frac13\log^3x\right]\dd x\notag\\
 &\hspace{20mm}=\bigl(\zeta(3)-\zeta(2)\bigr)
                                    (a^2-a+\tfrac16).
\end{align}
A compact table of multipliers is
\begin{center}
\begin{tabular}{cc l}
\toprule
$m$ & $q$ & $\kappa_{m,q}$\\
\midrule
1&1&$\zeta(2)$\\
1&2&$-2\zeta(3)$\\
1&3&$6\zeta(4)$\\
2&2&$\zeta(3)-\zeta(2)$\\
2&3&$3\zeta(3)-3\zeta(4)$\\
2&4&$12\zeta(5)-12\zeta(4)$\\
3&3&$\zeta(2)-3\zeta(3)+2\zeta(4)$\\
3&4&$-4\zeta(3)+12\zeta(4)-8\zeta(5)$\\
4&4&$-\zeta(2)+6\zeta(3)-11\zeta(4)+6\zeta(5)$\\
\bottomrule
\end{tabular}
\end{center}
The small size of these expressions is not a low-order accident: the
finite product $E_m$ fixes their length for every $q\ge m$.

\subsection{Reflection and moment rank}
For $0<a<1$, Bernoulli reflection gives
\begin{equation}\label{bsm:eq:reflection-moment}
 \int_0^\infty x^{m-1}
  \bigl(\Ccal_q(x,a)-(-1)^m\Ccal_q(x,1-a)\bigr)\dd x=0,
 \qquad q\ge m.
\end{equation}
Thus every odd $m\le q$ gives a centered vanishing moment at $a=1/2$.
For fixed $q\ge M$, the $M$ moment functions of $a$,
$\Jcal_q(1,a),\ldots,\Jcal_q(M,a)$, are linearly independent
polynomials. Indeed they are nonzero scalar multiples of
$B_1(a),\ldots,B_M(a)$, whose degrees are distinct. This is an exact
functional-rank statement, not a claim about the arithmetic independence
of special values.

\section{Exceptional moments and polylogarithm derivatives}\label{bsm:sec:exceptional}

\begin{theorem}[The Bernoulli-root moment]\label{bsm:thm:root}
Let $m\ge1$ and $a>0$ satisfy $B_m(a)=0$. Then the ordinary integral
\begin{equation}\label{bsm:eq:root}
 \boxed{\displaystyle
 \int_0^\infty x^{m-1}\Ccal_{m-1}(x,a)\dd x
          =(-1)^{m-1}\zeta^{[1]}(1-m,a)}
\end{equation}
converges absolutely.
\end{theorem}
\begin{proof}
Lemma~\ref{bsm:lem:tail} removes the leading $x^{-m}$ term, so the
integrand at infinity is $O(x^{-2})$. The integral is holomorphic
in $u$ in a neighborhood of $u=m$. Its value follows from
Theorem~\ref{bsm:thm:master} by continuation from $\Ree u<m$.
Here
\[
 \Qcal_{m-1}(m)=(m-1)!e_{m-1}^{(m)}=1,
 \qquad
 \zeta(1-u,a)=-(u-m)\zeta^{[1]}(1-m,a)+O((u-m)^2).
\]
Multiplication by $\pi\csc\pi u$ gives \eqref{bsm:eq:root}.
\end{proof}

This is a different mechanism from Theorem~\ref{bsm:thm:integer}. There,
$\Qcal_q(m)$ vanishes; here it does not, and the Hurwitz value vanishes
instead. Substituting $q=m-1$ into the finite sum
\eqref{bsm:eq:intro-kappa} would introduce $\zeta(1)$. Multiplying that
undefined expression by $B_m(a)=0$ is not a valid limiting argument.

\subsection{Centered and algebraic-shift examples}
For $m=1$, the only root of $B_1$ is $a=1/2$. The root formula becomes
\begin{equation}\label{bsm:eq:digamma-classical}
 \int_0^\infty\bigl(\log x-\psi(x+\tfrac12)\bigr)\dd x
               =\zeta^{[1]}(0,\tfrac12)=-\frac12\log2.
\end{equation}
This basic log-Gamma consequence is classical, not a novelty claim.
It is a useful normalization check on the higher-order family.

For $r\ge1$, both $B_{2r+1}(1)$ and $B_{2r+1}(1/2)$ vanish. Thus
\begin{align}
 \int_0^\infty x^{2r}\Ccal_{2r}(x,1)\dd x
 &=\frac{(-1)^r(2r)!}{2(2\pi)^{2r}}\zeta(2r+1),
 \label{bsm:eq:odd-one}\\
 \int_0^\infty x^{2r}\Ccal_{2r}(x,\tfrac12)\dd x
 &=\bigl(2^{-2r}-1\bigr)
       \frac{(-1)^r(2r)!}{2(2\pi)^{2r}}\zeta(2r+1).
 \label{bsm:eq:odd-half}
\end{align}
These use the standard zeta functional equation~\cite{DLMFReflection} and
$\zeta(s,1/2)=(2^s-1)\zeta(s)$. At $r=1$,
\begin{align}
 \int_0^\infty x^2\Ccal_2(x,1)\dd x
       &=-\frac{\zeta(3)}{4\pi^2},\label{bsm:eq:odd-first-one}\\
 \int_0^\infty x^2\Ccal_2(x,\tfrac12)\dd x
       &=\frac{3\zeta(3)}{16\pi^2}.\label{bsm:eq:odd-first-half}
\end{align}
At the two positive roots
$a_\pm=(1\pm1/\sqrt3)/2$ of $B_2$,
\begin{equation}\label{bsm:eq:B2-root}
 \int_0^\infty x\Ccal_1(x,a_\pm)\dd x
                  =-\zeta^{[1]}(-1,a_\pm).
\end{equation}
This is an exact reduction to a specified Hurwitz derivative. No
further elementary or cyclotomic evaluation at these algebraic shifts
is inferred.

\subsection{Every logarithmic moment at a Bernoulli root}

The extra convergence in Theorem~\ref{bsm:thm:root} persists after insertion
of any fixed power of $\log x$. It gives more than a single boundary
value.

\begin{corollary}[All exceptional logarithmic moments]\label{bsm:cor:root-logs}
Under the hypotheses of Theorem~\ref{bsm:thm:root}, for every $r\ge0$,
\begin{align}\label{bsm:eq:root-all-logs}
 &\int_0^\infty x^{m-1}\log^r x\,\Ccal_{m-1}(x,a)\dd x\notag\\
 &\quad=(-1)^{m-1}r![t^r]\left\{
   \frac{\pi t}{\sin\pi t}\Qcal_{m-1}(m+t)
   \sum_{j=0}^{r}\frac{(-1)^j\zeta^{[j+1]}(1-m,a)}{(j+1)!}t^j
                                      \right\}.
\end{align}
In particular, with
\begin{equation}\label{bsm:eq:alpha-m}
 \alpha_m=(m-1)!\sum_{j=0}^{m-2}(-1)^{m-j}e_j^{(m)}\zeta(m-j)
                    -H_{m-1},\qquad \alpha_1=0,
\end{equation}
the first logarithmic moment is
\begin{align}\label{bsm:eq:root-log-one}
 \int_0^\infty x^{m-1}\log x\,\Ccal_{m-1}(x,a)\dd x
   =(-1)^{m-1}\left\{
       \alpha_m\zeta^{[1]}(1-m,a)-\frac12\zeta^{[2]}(1-m,a)
                  \right\}.
\end{align}
\end{corollary}
\begin{proof}
The root condition makes $x^{m-1}\Ccal_{m-1}(x,a)=O(x^{-2})$.
Every displayed logarithmic integral therefore converges, and equals
the corresponding $u$-derivative of the Mellin transform at $u=m$.
Expand \eqref{bsm:eq:master} there, using
$\pi\csc\pi(m+t)=(-1)^m\pi/\sin\pi t$ and the vanishing Hurwitz
constant term. This proves \eqref{bsm:eq:root-all-logs}.
Since $\Qcal_{m-1}(m)=1$ and $\pi t/\sin\pi t$ has zero linear
coefficient, the first derivative gives \eqref{bsm:eq:root-log-one}
with $\alpha_m=\Qcal_{m-1}'(m)$.
To compute it explicitly, use \eqref{bsm:eq:Q-derivative} at degree
$m-1$. The finite rational recurrence term now contributes
$-H_{m-1}$, whereas the digamma Taylor series contributes the finite
zeta sum in \eqref{bsm:eq:alpha-m}.
\end{proof}

Two closed examples are
\begin{equation}\label{bsm:eq:root-log-digamma}
 \int_0^\infty\log x\bigl(\log x-\psi(x+\tfrac12)\bigr)\dd x
      =\frac14\log^22+\frac12\log2\log(2\pi)
\end{equation}
and
\begin{align}\label{bsm:eq:root-log-odd}
 &\int_0^\infty x^2\log x\,\Ccal_2(x,\tfrac12)\dd x\notag\\
 &\quad=\frac3{16\pi^2}\left[
 \zeta^{[1]}(3)+\zeta(3)\left(
 3\zeta(2)-2\zeta(3)-\gamma-\log(2\pi)+\frac{\log2}{3}
                                      \right)\right].
\end{align}
For the second formula, $\alpha_3=3\zeta(2)-2\zeta(3)-3/2$.
Differentiating the Riemann functional equation at $-2$ yields
\[
 \zeta^{[2]}(-2)=2\zeta^{[1]}(-2)
       \left(\log(2\pi)-\frac32+\gamma
                         -\frac{\zeta^{[1]}(3)}{\zeta(3)}\right).
\]
The half-shift duplication formula then gives
\eqref{bsm:eq:root-log-odd}. These formulas display explicitly how a
logarithmic moment introduces a positive-order zeta derivative.

\subsection{An explicit polylogarithmic form}
For $0<a<1$, define
\begin{equation}\label{bsm:eq:Ppoly}
 \mathcal P_\nu(a)=
 e^{-\ii\pi\nu/2}\Li_\nu(e^{2\pi\ii a})
 +e^{\ii\pi\nu/2}\Li_\nu(e^{-2\pi\ii a}).
\end{equation}
For $\Ree\nu>1$, all series and their finitely many order derivatives
converge absolutely. Hurwitz's Fourier formula is
\begin{equation}\label{bsm:eq:Hurwitz-polylog}
 \zeta(1-\nu,a)=\frac{\Gamma(\nu)}{(2\pi)^\nu}\mathcal P_\nu(a)
 \quad\cite[25.11.9]{DLMFHurwitz}.
\end{equation}
It uses the standard polylogarithm convention
$\Li_\nu(z)=\sum_{n\ge1}z^n/n^\nu$ in its initial disc and its
boundary continuation \cite{DLMFPolylog}.

\begin{corollary}[Polylogarithmic root identity]\label{bsm:cor:root-polylog}
If $m\ge2$, $0<a<1$, and $B_m(a)=0$, then
\begin{align}\label{bsm:eq:polylog-root}
 \int_0^\infty x^{m-1}\Ccal_{m-1}(x,a)\dd x
 =\frac{(-1)^m(m-1)!}{(2\pi)^m}
 \bigg\{&e^{-\ii\pi m/2}
   \left[\Li_m^{[1]}(e^{2\pi\ii a})
                 -\frac{\ii\pi}{2}\Li_m(e^{2\pi\ii a})\right]\notag\\
 +&e^{\ii\pi m/2}
   \left[\Li_m^{[1]}(e^{-2\pi\ii a})
                 +\frac{\ii\pi}{2}\Li_m(e^{-2\pi\ii a})\right]\bigg\}.
\end{align}
\end{corollary}
\begin{proof}
At the root, $\mathcal P_m(a)=0$ by
\eqref{bsm:eq:Hurwitz-polylog}. Differentiating that formula in $\nu$
therefore eliminates the derivative of its Gamma prefactor and gives
\[
 -\zeta^{[1]}(1-m,a)=\frac{(m-1)!}{(2\pi)^m}
                         \mathcal P_m^{[1]}(a).
\]
Use Theorem~\ref{bsm:thm:root} and differentiate
\eqref{bsm:eq:Ppoly}. The two terms in braces are conjugates for real $a$;
the displayed value is real.
\end{proof}

Away from the Bernoulli-root condition, the same differentiation
retains an additional $(\psi(m)-\log(2\pi))\mathcal P_m(a)$ term.
Dropping that term without its vanishing hypothesis would give a
wrong identity. For ordinary moments with $q\ge m\ge2$, one may
instead substitute
$B_m(a)=-m!\mathcal P_m(a)/(2\pi)^m$ directly into
\eqref{bsm:eq:integer-main}. This gives finite polylogarithm coordinates
at every rational shift without asserting their arithmetic independence.

\section{Finite shift filters: a complete moment criterion}\label{bsm:sec:filters}

Let $a_1,\ldots,a_J>0$ and $c_1,\ldots,c_J\in\C$. Write
\begin{equation}\label{bsm:eq:filter}
 \Lcal f=\sum_{j=1}^Jc_jf(a_j),\qquad
 Z_{\Lcal}(s)=\sum_{j=1}^Jc_j\zeta(s,a_j),\qquad
 \Ccal_q^{\Lcal}(x)=\sum_{j=1}^Jc_j\Ccal_q(x,a_j).
\end{equation}
The logarithmic term in $\Ccal_q^{\Lcal}$ vanishes only when
$\sum c_j=0$; until that condition is imposed it is retained.

\begin{theorem}[Necessary and sufficient cancellation]\label{bsm:thm:filter}
Let $m\ge1$, $q\ge0$. If $q\ge m$, the moment of
$\Ccal_q^{\Lcal}$ always converges absolutely and equals
$\kappa_{m,q}\Lcal B_m$. If $q<m$, the ordinary improper integral
\begin{equation}\label{bsm:eq:filter-moment}
 \int_0^\infty x^{m-1}\Ccal_q^{\Lcal}(x)\dd x
\end{equation}
converges if and only if
\begin{equation}\label{bsm:eq:filter-criterion}
 \Lcal B_k=0\qquad(k=q+1,\ldots,m).
\end{equation}
When it converges, it is absolutely convergent and
\begin{equation}\label{bsm:eq:filter-value}
 \boxed{\displaystyle
 \int_0^\infty x^{m-1}\Ccal_q^{\Lcal}(x)\dd x
   =(-1)^{m-1}q!e_q^{(m)}Z_{\Lcal}^{[1]}(1-m).}
\end{equation}
\end{theorem}

\begin{proof}
For $k\ge q+1$, $P_{k-1}^{(q)}(0)>0$: its polynomial has strictly
positive coefficients in all degrees from $0$ to $k-1$. Hence the
coefficients of $x^{-k}$ in Lemma~\ref{bsm:lem:tail} vanish for the filter
exactly when $\Lcal B_k=0$. If the first surviving power has $k\le m$,
the weighted integrand has a nonzero leading constant multiple of
$x^{m-1-k}$, whose integral diverges (logarithmically for $k=m$).
The lower powers of $x$ cannot cancel this divergence. This remains
true for complex weights because at least one real component has the
nonzero leading coefficient. If all powers through $k=m$ vanish,
the integrand is $O(x^{-2})$. The zero endpoint is integrable as before.
This proves necessity, sufficiency, and absolute convergence.

The filtered Mellin theorem initially holds on $0<\Ree u<q+1$ by
finite summation. Under \eqref{bsm:eq:filter-criterion} its left side
extends holomorphically through $u=m$. Since
$Z_{\Lcal}(1-m)=-\Lcal B_m/m=0$ and
$\Qcal_q(m)=q!e_q^{(m)}$, its removable limit is
\eqref{bsm:eq:filter-value}. The case $q\ge m$ is
Theorem~\ref{bsm:thm:integer} summed over the shifts.
\end{proof}

\subsection{Polynomial annihilators and minimal support}
Suppose $\Lcal$ annihilates every polynomial of degree at most $m$.
Equivalently,
\begin{equation}\label{bsm:eq:moment-filter}
 \sum_{j=1}^Jc_ja_j^k=0\qquad(k=0,\ldots,m).
\end{equation}
It then annihilates $B_0,\ldots,B_m$. Every balanced moment of order
$m$ exists and
\begin{equation}\label{bsm:eq:universal-filter}
 \int_0^\infty x^{m-1}\Ccal_q^{\Lcal}(x)\dd x=
 \begin{cases}
 (-1)^{m-1}q!e_q^{(m)}Z_{\Lcal}^{[1]}(1-m),&0\le q<m,\\
 0,&q\ge m.
 \end{cases}
\end{equation}
For distinct prescribed nodes, a nonzero filter satisfying
\eqref{bsm:eq:moment-filter} requires at least $m+2$ nodes. Indeed, the
Vandermonde matrix on any $m+1$ distinct nodes is invertible. With
$m+2$ nodes the nullspace is one-dimensional, so the normalized
annihilator is unique. This is a statement about full polynomial
annihilation, not about the weaker, potentially shorter list in
\eqref{bsm:eq:filter-criterion}.

For equally spaced nodes one may take the forward difference
\begin{equation}\label{bsm:eq:forward}
 \Delta_h^{m+1}f(a)=
 \sum_{j=0}^{m+1}(-1)^{m+1-j}\binom{m+1}{j}f(a+jh),
 \qquad a,h>0.
\end{equation}
Thus an entire infinite sequence of balanced moments vanishes after
one fixed finite difference.

\subsection{Concrete filtered identities}
Let
\begin{equation}\label{bsm:eq:half-filter}
 \Lcal f=f(\tfrac12)-2f(1)+f(\tfrac32).
\end{equation}
It kills constants and linear polynomials. At $m=1$, $q=0$,
\begin{equation}\label{bsm:eq:filtered-gamma0}
 \int_0^\infty\left[-\psi(x+\tfrac12)+2\psi(x+1)
                       -\psi(x+\tfrac32)\right]\dd x
             =\log\frac\pi2.
\end{equation}
At every $q\ge1$, the same filter applied to $\Ccal_q$ has integral
zero. The raw Stieltjes functions require the additional logarithmic
moment conversion developed next.

For the third difference with step $1/2$ based at $1/2$, put
\[
 \Lcal_3 f=-f(\tfrac12)+3f(1)-3f(\tfrac32)+f(2).
\]
The recurrence and duplication formulas give
\[
 Z_{\Lcal_3}(s)=4(2-2^s)\zeta(s)+3\cdot2^s-1.
\]
Its $m=2$ moments at $q=0$ and $q=1$ are both
\begin{equation}\label{bsm:eq:third-difference}
 \int_0^\infty x\Ccal_q^{\Lcal_3}(x)\dd x
       =-6\zeta^{[1]}(-1)-\frac53\log2\qquad(q=0,1),
\end{equation}
and every $q\ge2$ moment is zero. The identical first two values arise
from $0!e_0^{(2)}=1!e_1^{(2)}=1$, not from a presumed equality of
the two functions.

\section{All mixed logarithmic moments of filtered constants}\label{bsm:sec:mixed}

The finite-filter theorem concerns the balanced combinations. A stronger
local generating function recovers the separate filtered Stieltjes
constants with arbitrary powers of $\log x$.

\begin{theorem}[Complete mixed-moment formula]\label{bsm:thm:mixed}
Let $m\ge1$ and let $\Lcal$ satisfy \eqref{bsm:eq:moment-filter}.
For all integers $n,r\ge0$, the ordinary integral
\begin{equation}\label{bsm:eq:raw-moment}
 I_{n,r}^{(m)}(\Lcal)=
 \int_0^\infty x^{m-1}\log^r x
               \sum_j c_j\gamma_n(x+a_j)\dd x
\end{equation}
is absolutely convergent. Put $w=\eps-\eta$ and define the holomorphic
germ
\begin{align}\label{bsm:eq:K}
 \Kcal_m(\eps,\eta)
 &=\frac{\Gamma(1+\eta)\Gamma(1+w)}{\Gamma(1+\eps)}
       \prod_{j=1}^{m-1}\frac{1+\eta/j}{1-w/j},\notag\\
 M_m(\eps,\eta)
 &=(-1)^{m-1}\Kcal_m(\eps,\eta)
              \frac{Z_{\Lcal}(1-m+w)}w.
\end{align}
The quotient in the second line is removable because
$Z_{\Lcal}(1-m)=0$. Then
\begin{equation}\label{bsm:eq:mixed-coeff}
 \boxed{\displaystyle
 I_{n,r}^{(m)}(\Lcal)=
       (-1)^n n!r![\eps^n\eta^r]M_m(\eps,\eta).}
\end{equation}
Every fixed coefficient requires only
$Z_{\Lcal}^{[k]}(1-m)$ for $1\le k\le n+r+1$, ordinary zeta values,
and finite generalized harmonic numbers.
\end{theorem}

\begin{proof}
Because the filter kills every polynomial through degree $m$, the
first nonzero possible term in the large-$x$ expansion of
$\sum c_j\zeta(1+\eps,x+a_j)$ has index $k=m+1$. Its leading integral
term also cancels since $\sum c_j=0$. On a small compact neighborhood
of $(\eps,\eta)=(0,0)$, multiplication by $x^{m+\eta-1}$ therefore
gives an integrable bound of type $x^{-2+\delta}$, for some
$0<\delta<1$. Every parameter derivative adds only a fixed power of
$\log x$. At zero, all shifts are positive and
$x^{m+\eta-1}$ supplies the required integrability. The filtered
integral is consequently jointly holomorphic near the origin.

More explicitly, the same endpoint estimates show that the filtered
integral is jointly holomorphic on the connected tube
$\Ree(m+\eta)>0$, $\Ree(\eps-\eta)>-1$. The pole at
$\eps=0$ cancels because $\sum c_j=0$. This tube contains the
initial absolute beta domain $\Ree\eps>m+\Ree\eta$.
Summing \eqref{bsm:eq:raw-beta} there and continuing on this tube gives
\begin{equation}\label{bsm:eq:mixed-beta}
 \int_0^\infty x^{m+\eta-1}
         \sum_jc_j\zeta(1+\eps,x+a_j)\dd x
 =\frac{\Gamma(m+\eta)\Gamma(1-m+w)}{\Gamma(1+\eps)}
                                  Z_{\Lcal}(1-m+w).
\end{equation}
The possible Gamma poles in this tube occur at
$w=0,1,\ldots,m-1$. At each such point the corresponding
$Z_{\Lcal}(1-m+w)$ vanishes, since it is a multiple of one of the
annihilated Bernoulli polynomials $B_m,\ldots,B_1$.
Thus the complete right side is holomorphic as well, and the identity
theorem applies. Equivalently, finite Euler--Maclaurin subtraction
makes the same continuation explicit. No unspecified finite-part
constant is introduced.

The recurrence for Gamma factors gives
\[
 \Gamma(m+\eta)=(m-1)!\Gamma(1+\eta)
                       \prod_{j=1}^{m-1}(1+\eta/j),
\]
\[
 \Gamma(1-m+w)=\frac{(-1)^{m-1}\Gamma(1+w)}
                   {(m-1)!w\prod_{j=1}^{m-1}(1-w/j)}.
\]
Thus the right side of \eqref{bsm:eq:mixed-beta} is precisely
$M_m$. The coefficient of $\eps^n$ in the filtered Hurwitz family is
$(-1)^n\sum c_j\gamma_n(x+a_j)/n!$, and the coefficient of
$\eta^r$ in $x^\eta$ is $\log^r x/r!$. Extracting coefficients proves
\eqref{bsm:eq:mixed-coeff} and ordinary convergence of all its entries.
\end{proof}

\subsection{A finite harmonic and zeta coefficient compiler}
The logarithm of the elementary germ in \eqref{bsm:eq:K} is
\begin{align}\label{bsm:eq:K-log}
 \log\Kcal_m(\eps,\eta)
 ={}&\sum_{k\ge2}\frac{(-1)^k\zeta(k)}k
              \bigl(\eta^k+w^k-\eps^k\bigr)\notag\\
 &+\sum_{k\ge1}\frac{H_{m-1}^{(k)}}k
              \bigl((-1)^{k-1}\eta^k+w^k\bigr).
\end{align}
The Euler constant cancels because $\eta+w-\eps=0$.
For a requested total degree $d=n+r$, truncate this logarithm through
degree $d$, exponentiate formally through that degree, and multiply by
\begin{equation}\label{bsm:eq:Z-quotient}
 \frac{Z_{\Lcal}(1-m+w)}w
   =\sum_{j=0}^{d}\frac{Z_{\Lcal}^{[j+1]}(1-m)}{(j+1)!}w^j
                                      +O(w^{d+1}).
\end{equation}
No infinite coefficient operation is needed for a fixed identity.
The remaining Hurwitz derivatives are named analytic coordinates;
this algorithm alone does not prove they reduce to fewer constants.

The two simplest directions through this germ are
\begin{align}
 \Kcal_m(\eps,0)&=\prod_{j=1}^{m-1}(1-\eps/j)^{-1},
 \label{bsm:eq:K-vertical}\\
 \Kcal_m(\eps,\eps)&=\prod_{j=1}^{m-1}(1+\eps/j)=E_m(\eps).
 \label{bsm:eq:K-diagonal}
\end{align}
Thus raw order differentiation sees complete homogeneous harmonic
coefficients, while balanced diagonal differentiation sees elementary
ones. This explains both the finite degree and the infinite vanishing
family in \eqref{bsm:eq:universal-filter}.

\begin{corollary}[Raw moments without an extra logarithm]\label{bsm:cor:raw}
Let $h_j^{(m)}$ be defined by
$\prod_{k=1}^{m-1}(1-z/k)^{-1}=\sum_{j\ge0}h_j^{(m)}z^j$.
Under the hypotheses of Theorem~\ref{bsm:thm:mixed},
\begin{equation}\label{bsm:eq:raw-only}
 I_{n,0}^{(m)}(\Lcal)=(-1)^{n+m-1}n!
       \sum_{j=0}^{n}h_{n-j}^{(m)}
            \frac{Z_{\Lcal}^{[j+1]}(1-m)}{(j+1)!}.
\end{equation}
In particular,
\begin{equation}\label{bsm:eq:raw-m1}
 \int_0^\infty\sum_jc_j\gamma_n(x+a_j)\dd x
       =\frac{(-1)^n}{n+1}Z_{\Lcal}^{[n+1]}(0)
\end{equation}
when $\sum c_j=\sum c_ja_j=0$.
\end{corollary}
\begin{proof}
Use \eqref{bsm:eq:K-vertical} and \eqref{bsm:eq:Z-quotient}. For $m=1$ the
finite product is empty, so only $h_0^{(1)}=1$ survives.
\end{proof}

\subsection{A closed first-Stieltjes integral}
For the half-step filter \eqref{bsm:eq:half-filter},
\begin{equation}\label{bsm:eq:Z-half}
 Z_{\Lcal}(s)=2(2^s-2)\zeta(s)-2^s.
\end{equation}
The $n=1$ case of \eqref{bsm:eq:raw-m1} is therefore
\begin{align}\label{bsm:eq:gamma1-integral}
 &\int_0^\infty\bigl[
 \gamma_1(x+\tfrac12)-2\gamma_1(x+1)+\gamma_1(x+\tfrac32)
                         \bigr]\dd x\notag\\
 &\qquad=\gamma_1+\frac{\gamma^2}{2}-\frac{\pi^2}{24}
                     +\frac32\log^22-\frac12\log^2\pi.
\end{align}
Here $\gamma_1=\gamma_1(1)$. To check the constants, differentiating
\eqref{bsm:eq:Z-half} twice gives
\[
 -\tfrac12Z_{\Lcal}^{[2]}(0)
 =\zeta^{[2]}(0)-2\log2\,\zeta^{[1]}(0)+\log^22.
\]
Expansion of the standard Riemann functional equation at zero gives
\[
 \zeta^{[2]}(0)=\gamma_1+\frac{\gamma^2}{2}
                    -\frac{\pi^2}{24}-\frac12\log^2(2\pi),
 \qquad \zeta^{[1]}(0)=-\tfrac12\log(2\pi).
\]
Substitution proves \eqref{bsm:eq:gamma1-integral}.
Every higher $n$ similarly has the exact generating description
\eqref{bsm:eq:raw-m1}--\eqref{bsm:eq:Z-half}; expansion of the same functional
equation converts it to finitely many ordinary Stieltjes constants
and elementary constants. This is an all-order reduction mechanism,
not just the displayed first case.

For a genuinely mixed example with $m=1$, formula
\eqref{bsm:eq:K-log} starts with
$\log\Kcal_1=-\zeta(2)\eta(\eps-\eta)+O((|\eps|+|\eta|)^3)$.
Thus any filter killing constants and linear polynomials satisfies
\begin{equation}\label{bsm:eq:mixed-first}
 \int_0^\infty\log x\sum_jc_j\gamma_1(x+a_j)\dd x
       =\frac13Z_{\Lcal}^{[3]}(0)+\zeta(2)Z_{\Lcal}^{[1]}(0).
\end{equation}
The $\zeta(2)$ correction comes from the Gamma quotient. Omitting it
would confuse a raw moment with a balanced one.

\section{Shift derivatives and anchored primitives}\label{bsm:sec:calculus}

\begin{theorem}[Every shift derivative of the ordinary moments]\label{bsm:thm:derivatives}
For $q\ge m\ge1$, $a>0$, and $p\ge0$, differentiation in $a$ under
the integral is legitimate and
\begin{equation}\label{bsm:eq:derivative-moments}
 \int_0^\infty x^{m-1}\partial_a^p\Ccal_q(x,a)\dd x
 =\begin{cases}
 \kappa_{m,q}\dfrac{m!}{(m-p)!}B_{m-p}(a),&0\le p\le m,\\[5pt]
 0,&p>m.
 \end{cases}
\end{equation}
For $p\ge1$ the integrand has the explicit form
\begin{equation}\label{bsm:eq:derivative-kernel}
 \partial_a^p\Ccal_q(x,a)
 =(-1)^p\left.\partial_z^q
      \left[x^z(1+z)_p\zeta(p+1+z,x+a)\right]\right|_{z=0}.
\end{equation}
\end{theorem}
\begin{proof}
The identity $\partial_a^p\zeta(s,a)=(-1)^p(s)_p\zeta(s+p,a)$
gives \eqref{bsm:eq:derivative-kernel}. The Euler--Maclaurin expansion
used above is also normally uniform for $a$ in a small complex
neighborhood of any compact positive interval. Cauchy differentiation
therefore bounds every fixed $a$-derivative by the same integrable
endpoint type, possibly after taking extra expansion terms. One may
differentiate \eqref{bsm:eq:integer-main} under the integral. Finally
$B_m^{(p)}(a)=m!B_{m-p}(a)/(m-p)!$ for $p\le m$ and vanishes for
$p>m$.
\end{proof}

For example,
\begin{align}\label{bsm:eq:derivative-examples}
 \int_0^\infty\bigl[-\psi'(x+a)\log x-\gamma_1'(x+a)\bigr]\dd x
       &=\zeta(2),\notag\\
 \int_0^\infty\bigl[-\psi''(x+a)\log x-\gamma_1''(x+a)\bigr]\dd x
       &=0.
\end{align}
At the level of individual Stieltjes functions,
\begin{equation}\label{bsm:eq:Stieltjes-argument}
 \gamma_n^{(p)}(a)=(-1)^{n+p}n![z^n]
                         (1+z)_p\zeta(p+1+z,a),\qquad p\ge1.
\end{equation}
This is a finite sum of Hurwitz order derivatives with coefficients
from a harmonic polynomial. At $n=0$ it reduces to the usual
polygamma relation. For $n>0$ the order derivatives must not be
silently replaced by polygamma values alone.

\begin{theorem}[Anchored primitive and its moments]\label{bsm:thm:primitive}
For $a,b>0$, set
\begin{equation}\label{bsm:eq:primitive}
 \Acal_q(x;a,b)=
 -\frac1{q+1}\left.\partial_z^{q+1}
       \left\{x^z\bigl[\zeta(z,x+a)-\zeta(z,x+b)\bigr]\right\}
                                                    \right|_{z=0}.
\end{equation}
Then $\Acal_q(x;b,b)=0$ and $\partial_a\Acal_q=\Ccal_q(x,a)$.
Moreover, for $q\ge m\ge1$,
\begin{equation}\label{bsm:eq:primitive-moment}
 \int_0^\infty x^{m-1}\Acal_q(x;a,b)\dd x
   =\frac{\kappa_{m,q}}{m+1}
                  \bigl(B_{m+1}(a)-B_{m+1}(b)\bigr).
\end{equation}
All integrals in this statement are ordinary convergent integrals.
\end{theorem}
\begin{proof}
Differentiate \eqref{bsm:eq:primitive} in $a$ and use
$\partial_a\zeta(z,x+a)=-z\zeta(1+z,x+a)$. By
\eqref{bsm:eq:C-generator}, the coefficient of $z^{q+1}$ in
$z x^z\zeta(1+z,x+a)$ is $\Ccal_q(x,a)/q!$. The derivative of
order $q+1$, divided by $q+1$, is exactly $\Ccal_q$.
The anchoring is immediate. Thus
$\Acal_q(x;a,b)=\int_b^a\Ccal_q(x,t)\dd t$.
Uniform absolute convergence on the compact shift path permits Fubini
and \eqref{bsm:eq:integer-main}; integration of $B_m$ proves the moment.
\end{proof}

For direct use, the primitive is the finite expression
\begin{align}\label{bsm:eq:primitive-expanded}
 \Acal_q(x;a,b)= -\frac1{q+1}
 \sum_{j=0}^{q+1}\binom{q+1}{j}\log^{q+1-j}x\,
    \bigl[\zeta^{[j]}(0,x+a)-\zeta^{[j]}(0,x+b)\bigr].
\end{align}
The terms $j=0,1$ are elementary and log-Gamma coordinates. For
example,
\begin{equation}\label{bsm:eq:A0}
 \Acal_0(x;a,b)=(a-b)\log x-
                             \log\frac{\Gamma(x+a)}{\Gamma(x+b)}.
\end{equation}
Its derivative is the ordinary digamma kernel \eqref{bsm:eq:C0}.
Formula~\eqref{bsm:eq:primitive-expanded} specifies the integration
constant rather than hiding it in an unspecified antiderivative.

\section{Source audit, normalization, and computational evidence}\label{bsm:sec:audit}

\subsection{Scope of the source comparison}
The source inspection covered the canonical manuscript README,
the harmonic-generating and further-research sections
\cite{RepoHarmonic,RepoQuestions}, the Mellin continuation
\cite{RepoMellin}, selected related chapter text, and connector searches
for overlapping Stieltjes, Mellin, Bernoulli, and moment developments.
The repository was changing during retrieval. References therefore
record the actual snapshots or blob identifiers used rather than
pretending that all files came from one frozen checkout.

The incoming directory listing was inspected. It included archives
named \emph{Bernoulli Tails and Spectral Faces}, \emph{Higher Spectral
Defects}, \emph{Vertical Lattice Jets}, and other continuations. Their
binary archive interiors were not available in the working environment
for this session. Consequently we do not claim an exhaustive comparison
against those interiors or every preserved report. No theorem in this
article depends on an uninspected incoming claim.

The manuscript README correctly separates experimental evidence from
accepted proofs and retains $S_6$ and revised $S_8$ as conjectural
\cite{RepoREADME}. The harmonic-generating source correctly warns that
its finite-part constant does not determine all mixed harmonic
Stieltjes coefficients. We found no mathematical error in those
specific inspected statements.

\subsection{An actual editorial correction}
In the inspected README blob
\texttt{dc0fa26d0e5c2ede56df37c5640c7319681dfd25}, the paragraph beginning
with the complete $1{,}032$-input first-letter $X$ family, top-sector
rank $1{,}008$, and kernel dimension $24$ appears twice consecutively.
The proposed correction is to retain one copy, checking the current
file before applying a patch. This is an editorial duplication, not a
contradiction of the mathematical rank statement. The package contains
an integration note; no repository file was changed.

\subsection{Mathematical traps prevented by the present proofs}
There are three distinct cautions, none attributed to an existing
source without an actual occurrence.

First, the separate summands of \eqref{bsm:eq:C-intro} generally diverge
when integrated on $(0,\infty)$. The theorem concerns the combined
function. Second, the $q=m-1$ boundary must not be evaluated by
substituting into \eqref{bsm:eq:intro-kappa}: its last zeta value is
$\zeta(1)$, and the Bernoulli-root limit instead gives a zeta derivative.
Third, complete and elementary harmonic coefficients arise in
\emph{different} directions of \eqref{bsm:eq:K}; interchanging them destroys
both the finite-degree cancellation and the raw logarithmic corrections.

These warnings have explicit checks. At $q=0,m=1,a=1$, the tail is
$-1/(2x)$, so the proposed ordinary moment does not exist. At
$q=0,m=1,a=1/2$, the tail cancels and the correct value is
$-\log2/2$, not zero. At $m=1$, the mixed raw moment
\eqref{bsm:eq:mixed-first} contains the nonzero Gamma correction
$\zeta(2)Z_{\Lcal}^{[1]}(0)$.

\subsection{Independent exact checks}
The included \texttt{verify\_exact.py} performs exact rational and
symbolic comparisons. The independent constructions are integer
multiplication of Pochhammer polynomials versus symbolic products;
finite harmonic products versus a logarithmic Bell recurrence;
re-expansion of the Euler--Maclaurin terms versus Bernoulli polynomials;
exact linear systems for uniqueness; polynomial binomial identities
for the zeta multipliers; Vandermonde and finite-difference tests;
and elementary versus complete harmonic expansions of the two
mixed-germ restrictions. Two corruption controls deliberately reject
wrong coefficient patterns.

The delivered run passes \textbf{823 of 823 exact assertions}. These
finite checks verify the reported algebraic instances and the
implementation; they do not constitute proof-assistant formalization
of the all-index analytic statements. Those statements are supported
by the proofs in the preceding sections.

\subsection{Stable numerical evaluation}
Directly subtracting several large logarithmic Stieltjes terms at a
large $x$ is numerically wasteful. The verification code instead
uses an Euler--Maclaurin formula before taking the balanced derivative.
Let $A=x+a+N$ and $L=-\log(1+(a+N)/x)$. With $R$ correction terms,
\begin{align}\label{bsm:eq:EM-evaluator}
 \Ccal_q(x,a)\ \approx\ &
 \sum_{n=0}^{N-1}\frac{[-\log(1+(a+n)/x)]^q}{x+a+n}
       +\frac{L^{q+1}}{q+1}+\frac{L^q}{2A}\notag\\
 &+\sum_{r=1}^{R}\frac{B_{2r}}{(2r)!A^{2r}}
       \sum_{j=0}^{\min(q,2r-1)}\binom qj
                  P_{2r-1}^{(j)}(0)L^{q-j}.
\end{align}
The approximation sign is deliberate: the implemented remainder is
truncated, not interval-enclosed. The analytic identity underlying
this approximation has the usual Euler--Maclaurin remainder.

The integral from zero to a finite cutoff $T$ is evaluated by
tanh--sinh quadrature. On $(0,1)$ the substitution $x=t^4$ improves
resolution of logarithmic endpoints. The remaining tail is integrated
termwise from a finite asymptotic expansion \eqref{bsm:eq:C-tail}. For a
term $x^{-k}\log^r x$, its exact tail moment is
\begin{equation}\label{bsm:eq:log-tail}
 \int_T^\infty x^{u-k-1}\log^r x\dd x
 =T^{u-k}\sum_{j=0}^{r}\binom rj
               \frac{j!\log^{r-j}T}{(k-u)^{j+1}},\qquad\Ree u<k.
\end{equation}
Filtered coefficients are combined \emph{before} this step.
Nonintegrable powers are discarded only when their cancellation has
already been established analytically; floating-point smallness is
not used as a convergence theorem.

The main run uses $55$ working decimal digits, $N=56$, $R=40$,
$T=64$, and $64$ tail powers. Two selected identities are repeated
with $10$ additional digits and the enlarged parameters
$(N,R,T,K)=(72,50,80,80)$. The checks compare the integral side with
independently evaluated Mellin, Bernoulli, zeta-derivative, and mixed-germ
right sides. Pointwise comparisons with direct Stieltjes evaluations,
shift recurrences, and beta integrals provide additional controls.
The main run passes \textbf{34 of 34 diagnostics}, with maximum
scaled residual below $1.19\times10^{-53}$ at the stated settings.
Separate scripts add polylogarithm-root, primitive, averaging,
closed-constant, and exceptional logarithmic-moment checks. Exact per-case counts and residuals are
recorded in \texttt{results/numeric\_checks.json} and
\texttt{results/special\_checks.json}, together with
\texttt{results/log\_moment\_checks.json}. These are numerical diagnostics,
not rigorous accuracy certificates.

Direct near-integer order differentiation of the polylogarithm routine
also failed the demanded residual tolerance in three exploratory
checks. A $12$th-order centered stencil at $85$ working digits,
independently repeated at steps $10^{-5}$ and $10^{-6}$, removes that
numerical discrepancy. The superseded diagnostic is preserved in
\texttt{results/exploratory\_polylog\_derivative\_checks.json}; it is
not counted as a successful verification. The proof of the
polylogarithmic identity does not depend on either numerical route.

An initial low-resolution quadrature without the endpoint treatment
did not meet the accuracy target. It was not counted as verification.
This illustrates why high working precision alone does not certify
an improper integral. The delivered script records its method and
parameters explicitly.

\section{Further research questions}\label{bsm:sec:further}

\paragraph{1. Minimal filters for a prescribed collection of moments.}
Theorem~\ref{bsm:thm:filter} gives the exact Bernoulli constraints for one
pair $(m,q)$. For a finite set of pairs, determine the minimal support
and a canonical normalized solution for positive prescribed nodes.
The full polynomial annihilator is sufficient but can impose more
conditions than necessary. An exact classification of these smaller
Bernoulli evaluation matrices would sharpen the present finite-filter
construction without invoking arithmetic independence.

\paragraph{2. Contour transport for complex shifts.}
The local analytic arguments extend to $\Ree a>0$. Crossing a negative
real shift instead moves a Hurwitz singularity through the positive
integration axis. Develop contour-deformed analogues of the moment
formulas, with explicit contributions from residues and logarithmic
cuts. The objective is a branch-consistent identity, not continuation
of an unchanged real-axis integral across its singularities.

\paragraph{3. Function-level polylogarithmic identities.}
The present polylogarithmic evaluations come from the classical
Hurwitz functional equation. Seek independent functional relations
among $\mathcal P_\nu(a)$ and its order derivatives that cause entire
families of filtered moments to reduce to fewer coordinates. The
right target is an identity in parameters, not merely a small
numerical residual at a rational shift.

\paragraph{4. A colored or Lerch analogue of unique balancing.}
Replace Hurwitz zeta by a Lerch family with a nontrivial color. Its
order pole and large-shift expansion change, so the correct balancing
need not be \eqref{bsm:eq:C-intro}. Classify the homogeneous combinations
that cancel the first $q$ tails and determine whether their Mellin
moments reduce to Apostol--Bernoulli or other finite polynomial
coordinates. Rational colors should be compared with the finite
character filters already used in ProveIt.

\paragraph{5. Products of balanced kernels.}
The scalar Mellin transform is completely explicit. Products
$\Ccal_p(x,a)\Ccal_q(x,b)$ lead to Mellin convolutions and potentially
coupled multiple-zeta or multiple-Lerch coordinates. Determine a
first exact bilinear formula, including an honest function-class
statement. A product of two individually reduced transforms does not
by itself prove a reduction of the product integral.

\paragraph{6. Relation to genuinely harmonic Stieltjes data.}
The unresolved repository coefficients $R_{r,k}(a)$ belong to a
different family. Investigate whether the translation-variable
balancing used here can be combined with its moving-pole subtraction
to produce ordinary mixed moments. An answer should retain both
regulators and should not identify these coefficients with
$\gamma_k(a)$ by notation alone.

\paragraph{7. Higher anchored primitives and Barnes coordinates.}
Iterating Theorem~\ref{bsm:thm:primitive} introduces Hurwitz derivatives
at successively lower integers. Determine compact anchored formulas
using Barnes multiple Gamma functions, and compare normalizations
with the raw mixed-moment theorem. Constants at each anchor should
be explicit rather than absorbed into unrelated integration constants.

\paragraph{8. Certified computation and formal analytic interfaces.}
Turn the stable evaluator \eqref{bsm:eq:EM-evaluator} into an
outward-rounded algorithm with explicit remainder bounds. Separately,
formalize the cutoff Mellin continuation, locally normal parameter
differentiation, and finite-filter convergence criterion. These are
the reusable analytic steps; checking additional finite tables alone
would not formalize the all-order theorems.

\section{Conclusion}

The balanced generator $x^z\zeta(1+z,x+a)-1/z$ removes the order pole
and arranges the translation asymptotics into finite-degree
Pochhammer polynomials. This converts differentiation in $z$ into
exact tail cancellation. Its Mellin transform then exposes two
separate resonances: a zero of a harmonic Gamma coefficient gives
Bernoulli-polynomial moments, whereas a zero of the Bernoulli
polynomial gives an additional zeta-derivative moment. Finite filters
make the necessary cancellation conditions explicit and yield a
complete mixed logarithmic calculus.

The resulting identities are ordinary convergent integrals with
specified constants, not unnormalized finite parts. Their usefulness
for ProveIt lies in this combination of exact formulas, parameter
proofs, and a reproducible coefficient mechanism. They are offered
with a bounded source audit and a clear separation between analytic
proof, exact finite verification, numerical diagnostics, and questions
that remain unresolved.

\appendix
\section{Reproduction and integration map}\label{bsm:app:reproduction}

The package includes the complete source and compiled PDF in
\texttt{article/}, the mathematical routines and four verification
scripts in \texttt{code/}, and machine-readable evidence in
\texttt{results/}. The following commands run from the package root:
\begin{verbatim}
python code/verify_exact.py
python code/verify_numeric.py
python code/verify_special.py
python code/verify_log_moments.py
python build.py
\end{verbatim}
The first script requires SymPy and mpmath; the numerical script uses
mpmath. The build requires pdfLaTeX and the packages declared in the
preamble. The included \texttt{requirements.txt} specifies the versions
used for this delivery. The program does not require network access.

\begin{center}
\begin{tabular}{p{.35\textwidth}p{.51\textwidth}}
\toprule
Result & Principal proof dependency\\
\midrule
Unique balanced kernel & Uniform Hurwitz expansion; finite polynomial roots\\
Mellin theorem & Beta integral; explicit cutoff continuation; normal convergence\\
Integer moments & Gamma reflection; finite elementary harmonic coefficients\\
Exceptional moments & Bernoulli zero; removable Mellin singularity\\
Finite-filter criterion & Nonzero Pochhammer coefficients; leading-tail test\\
Mixed raw moments & Filtered beta transform; Gamma recurrence; joint holomorphy\\
Shift calculus & Hurwitz argument derivatives; compact normal convergence\\
\bottomrule
\end{tabular}
\end{center}

Suggested placement is a new continuation under the Stieltjes-correlation
Mellin material, with cross-references to the harmonic Gamma-coefficient
calculus. The report is self-contained and does not require promoting
any pending source theorem. Integration should retain its original
normalization, in particular the signs in \eqref{bsm:eq:C-intro}, the
factor $q!$ in \eqref{bsm:eq:intro-Q}, and the distinction between
\eqref{bsm:eq:integer-main} and \eqref{bsm:eq:root}. A short integration
fragment and source-status record are included. No automatic patch
or remote write is performed by the build or verification scripts.

\begin{thebibliography}{99}
\raggedright
\bibitem{DLMFHurwitz}
NIST Digital Library of Mathematical Functions,
\emph{\S25.11, Hurwitz Zeta Function}; especially equations 25.11.9,
25.11.12--14, 25.11.17--18, and the Euler--Maclaurin and asymptotic
subsections. Accessed 11 October 2026.
\url{https://dlmf.nist.gov/25.11}.

\bibitem{DLMFGamma}
NIST Digital Library of Mathematical Functions,
\emph{Chapter 5, Gamma Function}, including Gamma recurrence,
reflection, Taylor expansions, and polygamma relations.
\url{https://dlmf.nist.gov/5}.

\bibitem{DLMFPolylog}
NIST Digital Library of Mathematical Functions,
\emph{\S25.12, Polylogarithms}. Accessed 11 October 2026.
\url{https://dlmf.nist.gov/25.12}.

\bibitem{DLMFReflection}
NIST Digital Library of Mathematical Functions,
\emph{\S25.4, Reflection Formulas for the Riemann Zeta Function}.
Accessed 11 October 2026.
\url{https://dlmf.nist.gov/25.4}.

\bibitem{DLMFBernoulli}
NIST Digital Library of Mathematical Functions,
\emph{\S24.4, Basic Properties of Bernoulli and Euler Polynomials}.
Accessed 11 October 2026.
\url{https://dlmf.nist.gov/24.4}.

\bibitem{Zagier}
D. Zagier, \emph{The Mellin Transform and Related Analytic Techniques},
author-hosted exposition, especially the continuation by asymptotic
subtraction on pp.~1--3. Accessed 11 October 2026.
\url{https://people.mpim-bonn.mpg.de/zagier/files/tex/MellinTransform/fulltext.pdf}.

\bibitem{RepoREADME}
ProveIt Contributors, \emph{Polylogarithms and their Arithmetic Bridges},
canonical manuscript README, repository
\texttt{VladimirReshetnikov/ProveIt}, path
\texttt{Analysis/Polylogarithms/docs/manuscript/README.md}.
Inspected blob \texttt{dc0fa26d0e5c2ede56df37c5640c7319681dfd25}.

\bibitem{RepoHarmonic}
ProveIt Contributors, \emph{A Gamma coefficient calculus for shifted
harmonic sums}, file \texttt{sections/harmonic.tex} under
\texttt{stieltjes-harmonic-resonance/continuations/}\allowbreak
\texttt{stieltjes-harmonic-identities}.
Inspected blob \texttt{51108e8519d46af87dc3682956511fc525cecec3}.

\bibitem{RepoQuestions}
ProveIt Contributors, \emph{Further research questions and concrete
next targets}, file \texttt{sections/further\_research.tex} in the
same continuation as \cite{RepoHarmonic}.
Inspected blob \texttt{1051d5c4348e9b697712468e9c6b0c0f4d6373db}.

\bibitem{RepoMellin}
ProveIt Contributors, \emph{Gamma, polygamma, and Stieltjes identities},
file \texttt{sections/03-jets.tex} under
\texttt{stieltjes-correlation-calculus/continuations/}\allowbreak
\texttt{mellin-lerch-dilation}, retrieved at repository commit
\texttt{64c0a386f5eb0cfc7e1ebb751b46eaed20f3f84a}.
Full paths and the scope of the comparison are recorded in
\texttt{integration/SOURCE\_AUDIT.md}.
\end{thebibliography}
\end{document}
